\section{Finite algebraic regularity}\label{sec:finiteregularity}

We now pass from numerical spectral velocities to fixed algebraic
subspaces. Work first over $L$, where every norm is split. No common
apartment for the three boundary norms is assumed.

\subsection{Radial spectra and strict velocity gaps}

\begin{lemma}\label{lem:radialderivatives}
Fix a reference norm $o$, and let $h(u)$ be either vertex curve
of $c_*$. At every differentiability point of
$\lambda(u)=\sigma(o,h(u))$, its derivative is a permutation of
the intrinsic velocity list $v(u)$. For every $k$, the function
$u\mapsto S_k(o,h(u))$ is locally semiconvex.
\end{lemma}

\begin{proof}
We first give a uniform finite bound for the number of affine
pieces of the radial spectrum along any apartment geodesic
$\gamma$, of constant spectral velocity $v$. Apply the weighted
elimination in Lemma~\ref{lem:finitecommonbases} to the constant
norm $o$ and $\gamma$. A complete pivot sequence has at most
$(n!)^2$ possibilities, since there are at most $j^2$ choices at
residual size $j$. The validity set of each sequence is an interval,
defined by affine weight inequalities. Their endpoints subdivide
the parameter interval into at most $2(n!)^2+1$ open pieces.
On each, one fixed common basis works. Its relative weights have
derivatives that permute $v$. Ordering these affine weights adds
at most $n(n-1)/2$ cuts per piece. Thus a sufficient uniform bound is
\[
 N(n)=\bigl(2(n!)^2+1\bigr)\bigl(n(n-1)/2+1\bigr).
\]
This bound is independent of the basis, weights and interval length.

Fix a differentiability point $u_0$ of $\lambda$. Let $\gamma_h$
be the chord of $h$ on $[u_0-h,u_0+h]$, parametrized on that
interval, and put
\[
 v_h=\frac{\sigma(h(u_0-h),h(u_0+h))}{2h}.
\]
The quadratic chord estimate and
\eqref{eq:coordinatecontinuity} give $v_h\to v(u_0)$ and,
uniformly for $|s|\le h$,
\[
 \sigma(o,\gamma_h(u_0+s))
 =\lambda(u_0+s)+O(h^2)
 =\lambda(u_0)+s\lambda'(u_0)+o(h).
\]
The left side has an affine piece of length at least $2h/N(n)$.
Subtract the displayed equation at its endpoints and divide by
that length. For a permutation $P_h$ it follows that
$P_hv_h=\lambda'(u_0)+o(1)$. Finiteness of the permutation group
proves the assertion. This uses differentiability for one fixed
reference $o$; no common full-measure set for all reference norms
is needed.

The function $h\mapsto S_k(o,h)$ is geodesically convex by
Lemma~\ref{lem:spectralconvexity}, and is $k$-Lipschitz for
$d_\infty$. Comparison with the midpoint of the endpoint chord
therefore gives, locally uniformly,
\[
 S_k(o,h(u))\le\tfrac12\{S_k(o,h(u-h))+S_k(o,h(u+h))\}
                         +4kCh^2.
\]
Adding a sufficiently large quadratic makes this continuous
function midpoint convex, hence convex. This proves local
semiconvexity.
\end{proof}

\begin{lemma}\label{lem:topsaturation}
Suppose $v_k>v_{k+1}$ throughout an open interval $I$.
For every $s<t$ in $I$,
\begin{equation}\label{eq:topsaturation}
 S_k(h(s),h(t))=\int_s^t\sum_{i\le k}v_i(u)\,\dd u.
\end{equation}
\end{lemma}

\begin{proof}
We use a scalar consequence of semiconvexity. Let $F$ be semiconvex,
and suppose that almost everywhere its derivative belongs to a
finite collection of continuous functions $g_j$. Suppose their
maximum $M$ is locally Lipschitz and, on a compact interval, each
branch either equals $M$ identically or stays below $M-\delta$
for some $\delta>0$. If $F'_+(s)=M(s)$ at its left endpoint,
then $F'=M$ on that interval.

Indeed $G=F-\int_s^uM$ is semiconvex. Its right derivative is
right-continuous, has left limits, and obeys
$G'_-(u)\le G'_+(u)$. Approaching from the right through the
full-measure differentiability set, and taking a constant-index
subsequence, gives
\[
 G'_+(u)\in\{0\}\cup(-\infty,-\delta].
\]
At $s$ it is zero. Right-continuity makes the zero state persist
to the right. If the negative set had a first infimum after this
initial interval, its left limit there would be zero, so its right
derivative would be nonnegative and hence zero. Right-continuity
would again exclude negative values immediately afterwards,
a contradiction. Thus $G'_+=0$ throughout, and absolute
continuity proves the scalar assertion.

Apply this to $F(u)=S_k(h(s),h(u))$. By
Lemma~\ref{lem:radialderivatives}, its derivative almost everywhere
is a $k$-element subset sum of $v(u)$. The largest such sum is
$M_k=\sum_{i\le k}v_i$, and every other subset sum is smaller by
at least $v_k-v_{k+1}$. This has a positive minimum on any compact
subinterval of $I$. The forward secant limit gives
$F'_+(s)=M_k(s)$. The scalar assertion proves
\eqref{eq:topsaturation}.
\end{proof}

\begin{proposition}\label{prop:fixedgapsubspaces}
Suppose $v_k>v_{k+1}$ on an open interval containing $[l,r]$.
Choose a common orthogonal basis of $h(l),h(r)$, ordered by their
decreasing relative weights. Its first $k$ and last $n-k$ vectors
span subspaces $U,W$ for which
\begin{equation}\label{eq:fixedgapdecomposition}
 h(u)=h(u)|_U\oplus h(u)|_W\qquad(l\le u\le r).
\end{equation}
For every $l\le s<t\le r$, the relative spectrum on $U$ is the
first $k$ entries of the full spectrum. The block velocity lists
are therefore $(v_1,\ldots,v_k)$ and $(v_{k+1},\ldots,v_n)$.
Several strict cuts give one simultaneous fixed orthogonal block
decomposition, using the same endpoint basis.
\end{proposition}

\begin{proof}
Let $\tau_U$ be the wedge of the first $k$ endpoint basis vectors.
It realizes the top exterior operator norm between the endpoints.
For an intermediate $u$, exterior norm comparison gives
\[
 \begin{split}
 \log\frac{\det(h(u)|_U)(\tau_U)}{\det(h(l)|_U)(\tau_U)}
      &\le S_k(h(l),h(u)),\\
 \log\frac{\det(h(r)|_U)(\tau_U)}{\det(h(u)|_U)(\tau_U)}
      &\le S_k(h(u),h(r)).
 \end{split}
\]
The two left sides sum to $S_k(h(l),h(r))$. By
\eqref{eq:topsaturation}, the right sides have the same sum.
Both inequalities are equalities.

Apply the same argument to the wedge $\tau_W$ in the reverse
direction. The identity
\[
 S_{n-k}(h(t),h(s))
   =S_k(h(s),h(t))-d(h(t))+d(h(s))
\]
shows that this reverse partial sum is also saturated along the
chain. With the determinant generators from the endpoint basis,
we obtain
\[
 \begin{split}
 d(h(u)|_U)-d(h(l)|_U)&=S_k(h(l),h(u)),\\
 d(h(u)|_W)-d(h(l)|_W)&=d(h(u))-d(h(l))-S_k(h(l),h(u)).
 \end{split}
\]
The direct-sum norm $g_u=h(u)|_U\oplus h(u)|_W$ consequently has
the same determinant as $h(u)$, because the decomposition is
orthogonal at $l$. The ultrametric triangle inequality gives
$h(u)\le g_u$. Equal determinant and order force equality, proving
\eqref{eq:fixedgapdecomposition}.

Insert the chain $l,s,t,r$ in the exterior norm comparison. The
same saturation makes each inequality an equality. Thus the fixed
wedge $\tau_U$ realizes $S_k(h(s),h(t))$ for every pair.
Orthogonality makes the full spectrum the union of the two block
spectra, so the $U$ spectrum consists of its top $k$ entries,
including multiplicities. Secant limits identify the block
velocities. Their squared $d_2$ distances add to the full squared
distance, so the restricted curves are locally Lipschitz.
Finally, the same endpoint basis supplies nested prefix and suffix
subspaces for several cuts. Their intersections give the asserted
simultaneous decomposition.
\end{proof}

\subsection{Second-order scalar approximation}
A good time will mean a time $u_0$ where all $2n$ global velocity
entries are differentiable and where equal entries have equal
derivatives, also across the two vertices. These times have full
measure: a Lipschitz function has derivative zero almost everywhere
on its zero set, and there are only finitely many differences to
consider.

\begin{proposition}\label{prop:scalarjets}
At a good time, group all velocity entries by their value $p_g$.
Let $r_g$ be their common derivative, and set
$f_g(s)=p_gs+r_gs^2/2$, where $s=u-u_0$. There are fixed orthogonal
subspaces for these groups near $u_0$, and reference block norms
$h_{vg}$, such that
\begin{equation}\label{eq:scalarjets}
 \begin{split}
 \sup_{|s|\le h}d_2\bigl((c_*)_v(u_0+s),Q_v(u_0+s)\bigr)&=o(h^2),\\
 Q_v(u_0+s)&=\bigoplus_g e^{f_g(s)}h_{vg}.
 \end{split}
\end{equation}
In each group the intrinsic velocities are $f_g'(s)+o(h)$,
uniformly for $|s|\le h$. The reference pair $Q$ is ordered for
small $h$. Its reduced arrow is a direct sum over the common
values $p_g$ and is the actual reduced velocity decomposition at
$u_0$.
\end{proposition}

\begin{proof}
The gaps between the distinct values at $u_0$ persist nearby.
Proposition~\ref{prop:fixedgapsubspaces} gives fixed orthogonal
subspaces for these coarse groups, even when velocities within a
group separate at nearby times. Differentiability at $u_0$ gives,
for every entry in group $g$,
\[
 |v_{vgi}(u_0+s)-p_g-r_gs|\le\eta(h)|s|,
 \qquad |s|\le h,\qquad\eta(h)\to0.
\]
Normalize the restricted block curve by $e^{-f_g(s)}$. Its metric
speed is exactly the Euclidean norm of the block velocity entries
minus $f_g'(s)$, by the spectral secant formula for $d_2$.
The length bound for a locally Lipschitz metric curve therefore
bounds its distance from its value at $u_0$ by
$\sqrt{\dim V_{vg}}\eta(h)s^2/2$. Summing over the orthogonal
blocks proves \eqref{eq:scalarjets} and the velocity assertion.
The determinant logarithm of each block is $C^1$, with derivative
the sum of its continuous velocity entries.

Choose orthonormal $L$-bases for the reference block norms. In
these bases the original arrow matrix is integral. If a valuation-zero
entry connected distinct velocities $p_C\ne p_A$, the arrow
inequality for $Q$ would fail linearly on one side of $u_0$.
The $o(s^2)$ norm error in \eqref{eq:scalarjets} cannot repair
such a failure for the actual ordered curve. Thus every nonzero
cross-group entry has strictly positive valuation. There are only
finitely many entries, so these strict gaps persist for small $s$.
Within a common group the scalar functions $f_g$ at both vertices
agree, including their quadratic term. All its arrow inequalities
are those at $u_0$. Hence $Q$ is ordered. Its tangent agrees with
that of the actual curve, so its reduced summands are the actual
velocity blocks.
\end{proof}

\begin{lemma}\label{lem:cubicrepair}
At a good time, the reference $Q$ in
Proposition~\ref{prop:scalarjets} can be repaired at the endpoints
of $[u_0-h,u_0+h]$ to an ordered finite piecewise-apartment curve
$Q_h$ with the exact endpoints of $c_*$, such that $Q_h=Q$ on
$[u_0-h/2,u_0+h/2]$ and
\begin{equation}\label{eq:cubicrepair}
 E(Q_h)-\overline E(c_*|_{[u_0-h,u_0+h]})=o(h^3).
\end{equation}
\end{lemma}

\begin{proof}
Choose $\delta_h\downarrow0$ bounding the norm errors in
\eqref{eq:scalarjets} by $\delta_hh^2$ and the block velocity
errors by $\delta_hh$. Put $\ell_h=\sqrt{\delta_h}h<h/4$.
Keep $Q$ on the central interval and join it to the exact
$c_*$ endpoints by canonical bridges of length $\ell_h$.
The endpoints of each bridge are ordered and have the same fixed
block decomposition. The bridge respects that decomposition,
since its locus of orthogonal sums is geodesically convex.
Each pair of block endpoint norms has a common $L$-basis, so this
is a finite piecewise-apartment comparison curve.

After subtracting $f_g$, a bridge velocity differs from $f_g'$
by
\[
 O(\delta_hh^2/\ell_h+\ell_h)=O(\sqrt{\delta_h}h).
\]
The first term is the endpoint error divided by the bridge length;
the second compares the average quadratic slope to its derivative.
The original normalized block speed is $O(\delta_hh)$.
Throughout the interval $Q_h$ and $c_*$ are $O(\delta_hh^2)$
apart, including the quadratic chord error
$O(\ell_h^2)$ on the bridges. Both curves have a common strict
velocity margin for small $h$.

Use, separately in each block, the identity
\[
 \ell(z)=\ell(f_g')+\tan(f_g')(z-f_g')+D_\ell(z;f_g').
\]
On the strict margin the Bregman remainder is bounded by a
constant times $(z-f_g')^2$. Its integral for $c_*$ is
$O(\delta_h^2h^3)$; the bridges of $Q_h$ contribute
$O(\delta_h^{3/2}h^3)$, and its central part contributes zero.
The linear term is $R\tan(f_g')$ times the derivative of the
block determinant difference. Integrate it by parts. Its boundary
term vanishes separately in each block, because the repaired
endpoints are exact, and its integral is $O(\delta_hh^3)$.
The potential difference has the same bound. The original action
is its spectral integral by Proposition~\ref{prop:spectralaction};
the repaired curve has its ordinary differential action by
Lemma~\ref{lem:regularaction}. Combining these estimates proves
\eqref{eq:cubicrepair}.
\end{proof}

\subsection{Almost-everywhere semistability and straight directions}

\begin{proposition}\label{prop:aeharmonicity}
At almost every time the actual reduced velocity blocks are
semistable. A block has type $A^m$, $B^m$ or $C^m$, and its
position has derivative $-a$, $0$ or $a$, respectively, in the
coordinate $t=\sinh u$. In a mixed block the actual reduced arrow
is an isomorphism.
\end{proposition}

\begin{proof}
Work at a good time $u_0$ and use the quadratic reference and
repair above. For a group $g$, let $c_g,a_g$ be its source and
target dimensions, $d_g=c_g+a_g$, and let $R_0=R(u_0)$.
Choose a smooth function $\eta$ supported in $(-1/2,1/2)$.
On the central part of $Q_h$, vary all weights of this group by
$\varepsilon h^2\eta(s/h)$. Both signs are admissible for small
$h$: within the group the arrow ratios stay fixed, while the
strict cross-group gaps persist. The bridges are unchanged.

Expand the differential action, divide by $h^3$, and let
$h\to0$ with $\varepsilon$ fixed. The linear term is
\[
 \varepsilon\left\{
 -\frac{d_g}{2}\bigl(R_0\sec^2p_g\,r_g+R_0'\tan p_g\bigr)
       +qR_0(c_g-a_g)\right\}\int\eta.
\]
For the kinetic term, integrate $R\tan(f_g')$ against the
weight derivative; its constant part integrates to zero and its
first-order part gives the displayed derivative at $u_0$.
The remainder after division is $O(\varepsilon^2)+o(1)$,
uniformly for $\varepsilon$ in a fixed bounded interval.
Local minimality and \eqref{eq:cubicrepair}, followed by
$\varepsilon\to0$ from both sides, imply
\begin{equation}\label{eq:blockeuler}
 \frac{d_g}{2}\bigl(R_0\sec^2p_g\,r_g+R_0'\tan p_g\bigr)
       =qR_0(c_g-a_g).
\end{equation}

Let $S$ be a nonzero subrepresentation of the reduced block,
with dimensions $c_S,a_S$ and $d_S=c_S+a_S$. Choose integral
basis lifts of $S$ and complements at both vertices. Such changes
preserve the scalar reference $Q$ exactly. Reduced entries from
$S$ to its complement vanish, so their actual entries are zero
or have strictly positive valuation. Entries in the other
direction become more contracting when the weights on $S$
decrease. Thus, for $\eta\ge0$ and $\varepsilon\ge0$, the
following variation is ordered for small $h$: add
$(d_S/d_g)\varepsilon h^2\eta(s/h)$ to the whole group and
subtract $\varepsilon h^2\eta(s/h)$ on $S$.

All unperturbed group velocities are $f_g'$, and the total
weight shift, with dimension multiplicities, is zero. Its kinetic
first variation therefore vanishes pointwise. The potential
coefficient after division by $h^3$ is
\[
 \varepsilon qR_0
 \left\{\frac{d_S}{d_g}(c_g-a_g)-(c_S-a_S)\right\}\int\eta.
\]
Again the remainder is $O(\varepsilon^2)+o(1)$. If the bracket
were negative, choose $\varepsilon>0$ small enough that its
linear term dominates, then let $h\to0$. This contradicts local
minimality and \eqref{eq:cubicrepair}. Hence
\begin{equation}\label{eq:semistableslope}
 \frac{c_S-a_S}{c_S+a_S}\le\frac{c_g-a_g}{c_g+a_g}
 \quad\hbox{for every nonzero subrepresentation }S.
\end{equation}
For the charges in \eqref{eq:charges}, this is precisely the
semistability inequality: the phase increases strictly with the
ratio $(c-a)/(c+a)$.

Every finite-dimensional $A_2$ representation is the direct sum
of its kernel copies of $C$, image copies of $B$, and cokernel
copies of $A$, by choosing complements to the kernel and image.
Since their phases are distinct, semistability permits only one
of these types. Thus a shared velocity block is actually $B^m$,
with invertible reduced arrow, rather than only equal dimensions.
Finally, the derivative in $t$ of $x=R\tan v$ is, at the good
time,
\[
 \frac{\dd x}{\dd t}
 =\frac{R_0'\tan p_g+R_0\sec^2p_g\,r_g}{R_0}
 =a\frac{c_g-a_g}{c_g+a_g}.
\]
This gives $-a,0,a$ for the three types. Good times have full
measure, which proves the assertion.
\end{proof}

\subsection{Finiteness of the support}

\begin{proposition}\label{prop:finitesupport}
The spectral positions of $c_*$ form a finite balanced weighted
straight graph. It has at most $2n^2$ distinct finite event times
and at most $n^2$ essential geometric vertices.
\end{proposition}

\begin{proof}
Write $A_i(t)$ and $C_j(t)$ for the two ordered position lists.
They are locally Lipschitz. Proposition~\ref{prop:aeharmonicity}
gives, on one common set of full measure,
\[
 A_i'\in\{-a,0\},\qquad C_j'\in\{0,a\}.
\]
An $A$ position absent from the $C$ list has derivative $-a$,
and an unshared $C$ position has derivative $a$. At a shared
position the two multiplicities agree and the derivative is zero.

Each difference $C_j-A_i$ is nondecreasing. Its zero set is an
interval, possibly empty, a point, a ray or the whole line.
Moreover, if $s<t$ lie in that zero set, then
\[
 A_i(s)\ge A_i(t)=C_j(t)\ge C_j(s)=A_i(s),
\]
so both positions are constant between them. Let $D$ be the set
of finite endpoints of these $n^2$ zero sets. Then $|D|\le2n^2$.
On a component of $\RR\setminus D$, each opposite-type pair
either coincides throughout or never coincides. A paired branch
is constant; an unpaired $A$ branch has derivative $-a$ almost
everywhere and hence is exactly affine, and an unpaired $C$
branch is similarly affine of slope $a$.

These pieces define the graph with coincident positions combined
and their multiplicities retained. A constant branch is counted
once as a $B$ edge, with its common multiplicity. Parallel branches
of one free type do not intersect; a free branch cannot meet a
constant branch in the interior of a component, since that would
be an opposite-type contact. There are therefore no additional
crossings to insert there. The known tails supply its exterior
rays. This proves finiteness before any appeal to honeycomb
combinatorics.

At a vertex isolate its position from the other finitely many
positions at that time. Continuity implies that the same $A$
indices and the same $C$ indices enter and leave this neighborhood.
Writing $m_X^-$ and $m_X^+$ for the incoming and outgoing
multiplicities of type $X$, we obtain
\[
 m_A^-+m_B^-=m_A^++m_B^+,\qquad
 m_C^-+m_B^-=m_C^++m_B^+.
\]
The upward direction vectors
\[
 z_A=(-a/2,1/2),\qquad z_B=(0,1),\qquad z_C=(a/2,1/2)
\]
satisfy $z_B=z_A+z_C$. The preceding identities give their
weighted balance. There are no finite univalent vertices, and
the tails have degree $n$.

We may now apply the finite honeycomb construction of
\cite[\S3.1, Lemmas~4 and 6--8, and Theorem~1]{KT99}.
Its dual is a subdivision of the lattice triangle of side $n$.
The two-dimensional cells correspond to essential geometric
vertices, with each cell a union of elementary triangles.
The entire triangle has $n^2$ elementary triangles. Every cell
contains at least one, proving the bound $n^2$.
Degenerate or crossing vertices correspond to larger cells;
artificial bivalent subdivisions are removed. This last count
bounds vertices, whereas the earlier direct estimate bounds
only their distinct time coordinates.
\end{proof}

\subsection{From finite support to finite apartments}

\begin{proposition}\label{prop:finitehahnregularity}
The Hahn minimum has a finite fixed-apartment presentation. On each
closed finite piece its weights extend real analytically beyond the
endpoints. On every open piece its actual reduced velocity blocks
are semistable and have the straight directions prescribed by their
charges.
\end{proposition}

\begin{proof}
Choose finite cutoffs in the known algebraic tails and subdivide
between them at the support event times. On an open piece $(b,c)$,
the ordered position branches are affine in $t$. Distinct branches
do not intersect there, and coincident branches agree throughout.
Thus the velocity groups have constant multiplicities and functions
\begin{equation}\label{eq:analyticvelocities}
 p_g(u)=\arctan(\alpha_g\tanh u+\beta_g\sech u),
 \qquad \alpha_g\in\{-a,0,a\}.
\end{equation}
Each extends analytically beyond the closed piece. The gaps
between distinct groups are positive in its interior but may
vanish at either endpoint.

For each cut $k$ between distinct groups,
\eqref{eq:topsaturation} holds for all $b<s<t<c$.
Continuity of the norms, spectra and velocities extends it to
$b\le s<t\le c$. Choose one common endpoint basis, ordered by
relative weights. The extremal-wedge proof of
Proposition~\ref{prop:fixedgapsubspaces} uses only these saturated
identities on closed chains. It therefore gives one simultaneous
orthogonal block decomposition on the whole closed interval,
even when a velocity gap vanishes at an endpoint. No lower bound
for that endpoint gap is needed. If there is only one group,
take the entire space as its block.

In block $g$ all velocities equal $p_g$. Fix an interior time
$u_*$, set $f_g(u)=\int_{u_*}^u p_g(s)\,\dd s$, and let
$h_g$ be the block norm at $u_*$. The normalized block curve
$e^{-f_g(u)}c_*(u)|_{V_g}$ has metric speed zero. Its local
Lipschitz continuity and the metric length bound make it constant.
Hence
\[
 c_*(u)|_{V_g}=e^{f_g(u)}h_g
\]
throughout the interior and, by continuity, at both endpoints.
Choose an orthogonal basis for each split norm $h_g$. Their union
is one fixed apartment on the entire closed piece, and its weights
are $f_g$ plus constant offsets. This proves the claimed analytic
extension without taking a limit of bases on an exhaustion.

In these fixed bases, equal-velocity coordinates have constant
relative weights, so their normalized reduced arrow matrix is
constant. A nonzero entry between distinct groups has strictly
positive norm gap at every interior time: an interior zero of
its nonnegative smooth gap would have zero derivative and would
force the two velocities to agree. Thus these fixed groups are
the actual reduced blocks everywhere on the open piece.
Their almost-everywhere semistability from
Proposition~\ref{prop:aeharmonicity}, together with the constant
reduced matrices, proves semistability throughout. Adjoin the
finitely many tail apartments to complete the presentation.
\end{proof}

\subsection{Finite descent and identification of the original minimum}
The preceding presentation is finite, but its matrices lie in $L$.
The following finite realization preserves more than spectra: it
retains the actual reduced arrows and the full endpoint norms.

\begin{lemma}\label{lem:finitecurvetransfer}
Let an ordered continuous $L$-norm pair on a compact interval have
finitely many fixed-apartment pieces, with prescribed real $C^1$
weights and strict velocity margins. Suppose its endpoint norms
are extensions of specified split $K$-norms. There is a continuous
ordered $K$-norm pair with the same partition and weights, exactly
those full endpoints, the same determinant functions, the same
normalized reduced arrow matrices on all pieces, and the same
normalized reduced transition matrices at their joins. In
particular its differential action is unchanged.
\end{lemma}

\begin{proof}
Use the same ambient $K$-frame at both quiver vertices. Let
$P_{vj}$ be the finitely many original apartment basis matrices.
Choose orthonormal $K$-frames $F_{v,-},F_{v,+}$ for the prescribed
outer endpoint norms. Introduce matrix variables $Q_{vj}$ and
$R_{vj}$ with $Q_{vj}R_{vj}=R_{vj}Q_{vj}=I$. Consider the finite
list consisting of the entries of
\[
 \begin{gathered}
 R_{Aj}Q_{Cj},\qquad
 R_{vj}Q_{v,j+1},\quad R_{v,j+1}Q_{vj},\\
 R_{v,\mathrm{first}}F_{v,-},\quad
 F_{v,-}^{-1}Q_{v,\mathrm{first}},\quad
 R_{v,\mathrm{last}}F_{v,+},\quad
 F_{v,+}^{-1}Q_{v,\mathrm{last}},
 \end{gathered}
\]
and of the determinants $\det Q_{vj}$. Evaluate them at the
original matrices $P_{vj}$ and their inverses. Preserve every
zero by an exact equation. For a nonzero value $x$, put
$\rho=\nu(x)$ and $a_x=\operatorname{res}(T^{-\rho}x)$, and
impose the valued-field condition
\begin{equation}\label{eq:leadingresiduecondition}
 T^{-\rho}x-a_x\in\mm.
\end{equation}
All parameters lie in $K$, including the finitely many real
monomials and complex residues. The original $L$-matrices satisfy
this finite formula, so elementarity gives invertible $K$-matrices
satisfying it. It preserves each entry valuation and leading
residue, and each determinant valuation.

For a matrix $M$ from a source basis of weights $w_j$ to a target
basis of weights $z_i$, nonexpansiveness is the finite collection
of inequalities
\begin{equation}\label{eq:matrixorder}
 \nu(M_{ij})+w_j-z_i\ge0
\end{equation}
for its nonzero entries. The valuations and weight functions are
unchanged. Thus the transferred arrow is nonexpansive at every
time. Both directions of each transition were included, so their
inequalities prove equality of the two norm descriptions at each
join. The same two-way argument with the fixed endpoint frames
proves equality with the full prescribed endpoint norms, on every
vector. It is stronger than preservation of their spectra.

A normalized arrow entry is $T^{w_j-z_i}M_{ij}$. If its valuation
is positive, its reduction is zero; if it is zero, its reduction
is the leading residue preserved in
\eqref{eq:leadingresiduecondition}. Hence all the reduced matrices
agree in the indexed bases, both on pieces and at joins.
Equation~\eqref{eq:detnorm} shows that the determinant functions
are unchanged. The weights and their derivatives are also
unchanged, so the differential actions agree. No time variable
or infinite family of norm inequalities was quantified: the finite
entry data imply them for the fixed weight functions.
\end{proof}

\begin{theorem}\label{thm:finitealgebraicminimum}
For every full split boundary triple, the original $K$-minimum
$c_*$ has a finite regular harmonic presentation as in
Definition~\ref{def:regularharmonic}. It has the finite weighted
support of Proposition~\ref{prop:finitesupport}, satisfies the
normalized ends, and obeys the stronger $O(h^2)$ midpoint estimate
at every join.
\end{theorem}

\begin{proof}
Choose $l<0<r$ inside the established fixed $K$-apartment tails.
Proposition~\ref{prop:finitehahnregularity} gives a finite Hahn
presentation on $[l,r]$. Apply
Lemma~\ref{lem:finitecurvetransfer} to its exact split $K$
endpoint norms, and append the actual known $K$ tails outside
this interval. Call the resulting curve $z_K$.
It is ordered and has the same weights, determinant functions
and reduced labels as the Hahn minimum in the middle. Its tails
are unchanged. We show that it is the actual $K$-minimum,
rather than merely a curve with the same numerical data.

For either finite presentation, its differential $J$ is finite.
On the negative tail the gap is zero and the velocities are
$O(e^u)$, so the integrand is $O(e^u)$. On the positive tail the
velocity difference from $\epsilon_v\beta$ is $O(e^{-u})$,
so the Bregman integrand is $O(e^{-u})$ as well. The norm errors
from the reference tails have the same rates.

Cut the actual tails at $\pm V$, keep the middle curve, and
join to the exact reference tails on intervals of length one.
At the negative end the common norms can be bridged together,
so their gap stays zero and their velocities are $O(e^{-V})$.
At the positive end both endpoint pairs are ordered; each
vertex bridge is in its fixed tail apartment. Relative to the
reference velocity its error is $O(e^{-V})$, since the endpoint
norm errors have that size and $\beta'=O(e^{-2V})$.
The bridge and removed tail costs are $O(e^{-V})$. These curves
belong to $\calQ$ and converge uniformly to the given finite
presentation. In particular,
\[
 J_*^L\le J_{\mathrm{diff}}(c_*^L).
\]

For the reverse inequality, take a global polygonal minimizing
sequence converging uniformly to $c_*^L$. On each fixed
$[-V,V]$, the partition lower bound gives
\[
 E_{\mathrm{diff},[-V,V]}(c_*^L)
       \le\liminf_m E_{[-V,V]}(p_m).
\]
The endpoint determinant terms converge. Subtracting the finite
correction in \eqref{eq:renormalization} gives the same lower
bound for truncated $J$. Since its integrands are nonnegative,
$J_{\mathrm{diff},[-V,V]}(c_*^L)\le J_*^L$.
Monotone convergence as $V\to\infty$ proves
$J_{\mathrm{diff}}(c_*^L)=J_*^L$.

Scalar extension equates the two field infima. The transferred
compact actions and the identical tails now give
\[
 J_{\mathrm{diff}}(z_K)=J_{\mathrm{diff}}(c_*^L)=J_*^K.
\]
The same exact-tail recovery over $K$ places $z_K$ in $\calY_K$
and gives $\overline J(z_K)\le J_*^K$. Uniqueness in
Theorem~\ref{thm:relaxedminimum} forces $z_K=c_*^K$ as actual
norm curves. Thus the original minimum has the finite $K$
presentation and actual semistable labels asserted here.
Its midpoint estimate is Proposition~\ref{prop:quadraticchord};
it does not depend on the transferred coordinates. The already
proved tails give precisely the full normalized end limits.
\end{proof}

The construction has supplied algebraic edge labels on a finite
support. To finish the path correspondence in
Theorem~\ref{thm:main}, we must construct the local vertex
filtrations and prove that every finite regular path with
first-order matching is this minimum. Those are distinct uses
of the matching condition, addressed next.
